\documentclass[8pt, a4paper]{article}

% ── Packages ─────────────────────────────────────────────────────────────────
\usepackage[a4paper, top=2.5cm, bottom=2.5cm, left=2.8cm, right=2.8cm]{geometry}
\usepackage[T1]{fontenc}
\usepackage[utf8]{inputenc}
\usepackage{lmodern}
\usepackage{microtype}
\usepackage{parskip}          % space between paragraphs, no indent
\usepackage{titlesec}         % section heading styles
\usepackage{enumitem}         % list customisation
\usepackage{booktabs}         % professional tables
\usepackage{tabularx}         % auto-width table columns
\usepackage{array}            % extra column formatting
\usepackage{xcolor}           % colours
\usepackage{mdframed}         % framed boxes
\usepackage{fancyhdr}         % header / footer
\usepackage{amsmath}          % math environments (cases)
\usepackage{hyperref}         % clickable links (load last)

% ── Colours ───────────────────────────────────────────────────────────────────
\definecolor{accentblue}{RGB}{30, 90, 160}
\definecolor{rulegray}{RGB}{190, 195, 205}
\definecolor{boxbg}{RGB}{237, 242, 250}

% ── Hyperref setup ────────────────────────────────────────────────────────────
\hypersetup{
  colorlinks = true,
  linkcolor  = accentblue,
  urlcolor   = accentblue,
}

% ── Section heading styles ────────────────────────────────────────────────────
\titleformat{\section}
  {\large\bfseries\color{accentblue}}
  {\thesection.}{0.6em}{}
  [\vspace{-4pt}\textcolor{rulegray}{\rule{\linewidth}{0.6pt}}\vspace{2pt}]

\titlespacing{\section}{0pt}{18pt}{6pt}

% ── Concept box ───────────────────────────────────────────────────────────────
\newmdenv[
  backgroundcolor   = boxbg,
  linecolor         = accentblue,
  linewidth         = 1.2pt,
  leftline          = true,
  rightline         = false,
  topline           = false,
  bottomline        = false,
  innerleftmargin   = 10pt,
  innertopmargin    = 6pt,
  innerbottommargin = 6pt,
  skipabove         = 6pt,
  skipbelow         = 6pt,
]{conceptbox}

% ── List styling ──────────────────────────────────────────────────────────────
\setlist[itemize]{topsep=3pt, itemsep=2pt, parsep=0pt, leftmargin=1.6em}
\setlist[enumerate]{topsep=3pt, itemsep=2pt, parsep=0pt, leftmargin=1.8em}

% ── Header / Footer ───────────────────────────────────────────────────────────
\pagestyle{fancy}
\fancyhf{}
\renewcommand{\headrulewidth}{0.4pt}
\fancyhead[L]{\small\color{accentblue}\textbf{ELEC3609 / ELEC9609}}
\fancyhead[R]{\small\color{gray}Week 3 --- Web Servers \& Application Architecture}
\fancyfoot[C]{\small\thepage}

% =============================================================================
\begin{document}

% ── Title Block ───────────────────────────────────────────────────────────────
\begin{center}
  {\LARGE\bfseries\color{accentblue} ELEC3609 / ELEC9609: Internet Software Platforms}\\[6pt]
  {\large Week 3 Study Guide: Web Servers \& Application Architecture}
\end{center}

\vspace{12pt}
\textcolor{rulegray}{\rule{\linewidth}{1pt}}
\vspace{6pt}

% =============================================================================
\section{Three-Tier Web Architecture \& Stacks}

Modern web platforms are organized into tiers to separate concerns, improve security, and simplify scalability.

\begin{conceptbox}
\textbf{The Three Architectural Layers}
\begin{enumerate}
  \item \textbf{Web Server Layer:} Handles public incoming HTTP (Port 80) and HTTPS (Port 443) traffic. Serves static assets (HTML, CSS, JavaScript, Images) directly. Acts as a reverse proxy forwarding dynamic requests to the Application Server.
  \item \textbf{Application Server Layer:} Executes core business logic and dynamic request processing. Communicates via HTTP, WSGI, or custom TCP protocols. Tightly integrated with software frameworks (e.g., Django, Flask, Express).
  \item \textbf{Data Storage Layer:} Persists stateful data (Relational DBs, NoSQL, Object Storage, Caches). Placed in private networks behind firewalls (isolated from public Internet).
\end{enumerate}
\end{conceptbox}

\begin{conceptbox}
\textbf{Classic Web Stacks}
\begin{itemize}
  \item \textbf{LAMP:} Linux (OS) + Apache (Web Server) + MySQL (Database) + PHP (Language)
  \item \textbf{LEMP:} Linux (OS) + Nginx (Web Server) + MySQL (Database) + PHP (Language)
  \begin{itemize}
    \item Note: ``E'' in LEMP comes from ``Engine-X'' (Nginx).
  \end{itemize}
\end{itemize}
\end{conceptbox}

\begin{conceptbox}
\textbf{Web Server vs.\ Application Server}

\smallskip
\begin{tabularx}{\linewidth}{l X X}
  \toprule
  \textbf{Feature} & \textbf{Web Server} & \textbf{Application Server} \\
  \midrule
  Primary Function    & Serves static content \& HTTP(S)                         & Executes dynamic business logic \\
  Protocols Supported & Exclusively HTTP / HTTPS                                 & HTTP, TCP, WSGI, gRPC, RPC, etc. \\
  Speed / Overhead    & Extremely high throughput, low overhead for static files & Higher processing load per request due to code execution \\
  Examples            & Nginx, Apache HTTP Server, IIS                           & Gunicorn, uWSGI, Tomcat, Node.js \\
  \bottomrule
\end{tabularx}
\end{conceptbox}

% =============================================================================
\section{Web Server Engines Comparison (IIS, Apache, Nginx)}

Choosing the correct web server depends on scale requirements and OS ecosystem.

\begin{conceptbox}
\textbf{1.\ Microsoft IIS (Internet Information Services)}
\begin{itemize}
  \item \textbf{Launched:} 1996 (Microsoft)
  \item \textbf{Architecture:} Proprietary Windows web server.
  \item \textbf{Best For:} Native ASP.NET applications and enterprise Active Directory (AD) integration.
  \item \textbf{Drawbacks:} Historically prone to security vulnerabilities; bound to Windows environment.
\end{itemize}
\end{conceptbox}

\begin{conceptbox}
\textbf{2.\ Apache HTTP Server}
\begin{itemize}
  \item \textbf{Launched:} 1995 (Apache Software Foundation)
  \item \textbf{Architecture:} Process/Thread-per-request model (\texttt{mpm\_prefork}, \texttt{mpm\_worker}).
  \item \textbf{Features:} Highly modular using dynamically loaded modules (\texttt{mod\_rewrite}, \texttt{mod\_ssl}).
  \item \textbf{Scale Characteristics:} Flexible, but higher memory footprint under massive concurrency.
\end{itemize}
\end{conceptbox}

\begin{conceptbox}
\textbf{3.\ Nginx (``Engine-X'')}
\begin{itemize}
  \item \textbf{Launched:} 2004 (Igor Sysoev)
  \item \textbf{Architecture:} Asynchronous, event-driven worker process model.
  \item \textbf{Features:} Acts as web server, reverse proxy, and layer 7 load balancer.
  \item \textbf{Scale Characteristics:} Consistent low CPU and memory consumption under heavy traffic; designed specifically to solve the C10K problem (10,000 concurrent connections).
\end{itemize}
\end{conceptbox}

% =============================================================================
\section{Programming Languages \& Frameworks (Django MVT Model)}

\begin{conceptbox}
\textbf{Web Server Gateway Interface (WSGI)}
\begin{itemize}
  \item WSGI is the standard Python specification (PEP 3333) bridging Web Servers (Nginx/Apache) with Python Web Applications (Django/Flask).
  \item Nginx/Apache cannot run Python code natively; WSGI servers (Gunicorn/uWSGI) translate HTTP requests into Python function calls.
\end{itemize}
\end{conceptbox}

\begin{conceptbox}
\textbf{Django MVT Architecture vs.\ Traditional MVC}

Django uses the MVT (Model-View-Template) pattern, which maps directly to classic MVC:

\smallskip
\begin{tabularx}{\linewidth}{l X}
  \toprule
  \textbf{Django Component} & \textbf{Role \& Responsibilities} \\
  \midrule
  Model (M)    & Encapsulates data structure and database operations (ORM). Maps directly to `Model' in classic MVC. \\
  View (V)     & Receives HTTP requests, executes business logic, fetches data from Models, and returns HTTP responses. Maps to `Controller' in classic MVC! \\
  Template (T) & Presentation layer using HTML \& Django Template Language. Maps to `View' in classic MVC! \\
  \bottomrule
\end{tabularx}
\end{conceptbox}

\begin{conceptbox}
\textbf{MVT Request-Response Lifecycle Flow}

\medskip
Client HTTP Request $\to$ Web Server $\to$ WSGI (Gunicorn) $\to$ Django URL Router (\texttt{urls.py})\\
\hspace*{1.5em}$\to$ Django View (\texttt{views.py}) $\leftrightarrow$ Model / DB (\texttt{models.py})\\
\hspace*{1.5em}$\to$ View populates Template (HTML/DTL) $\to$ HTTP Response sent to Client.
\end{conceptbox}

% =============================================================================
\section{Data Storage Landscape: Relational, NoSQL \& Cloud Storage}

\begin{conceptbox}
\textbf{Storage Types \& Form Factors}
\begin{itemize}
  \item \textbf{In-Memory Storage:} Ultra-fast volatile RAM storage (e.g., Redis, Memcached).
  \item \textbf{Block Store:} Raw unformatted volumes attached to VMs (e.g., AWS EBS, local HDDs/SSDs).
  \item \textbf{File Store:} Shared network file systems (e.g., NFS, Samba, CephFS, HDFS).
  \item \textbf{Object Store:} Highly durable key-value store for unstructured files/media (e.g., AWS S3, GCP Cloud Storage).
  \item \textbf{SAN vs.\ NAS:}
  \begin{itemize}
    \item \textbf{SAN (Storage Area Network):} Block-level storage over Fibre Channel or iSCSI.
    \item \textbf{NAS (Network Attached Storage):} File-level storage over local network protocols (NFS/SMB).
  \end{itemize}
\end{itemize}
\end{conceptbox}

\begin{conceptbox}
\textbf{Relational Databases (RDBMS)}

Organizes data into structured tables with rows (instances) and columns (attributes). Enforces schema rigidity (schema updates usually require planned maintenance/downtime).

\medskip
\textbf{1.\ MySQL / MariaDB:}
\begin{itemize}
  \item Most widely deployed open-source RDBMS.
  \item Non-standard SQL gotcha: \texttt{||} operates as logical \texttt{OR} (whereas standard SQL uses \texttt{||} for string concatenation).
\end{itemize}

\textbf{2.\ PostgreSQL:}
\begin{itemize}
  \item Strictly ANSI/ISO SQL compliant with complete ACID transaction guarantees.
  \item Outstanding spatial features (PostGIS) and advanced data types (JSONB).
\end{itemize}

\textbf{3.\ Microsoft SQL Server (MSSQL):}
\begin{itemize}
  \item Deep integration with MS Enterprise BI stack (SSIS, SSAS, SSRS).
  \item High licensing cost overhead; requires careful DBA planning.
\end{itemize}
\end{conceptbox}

\begin{conceptbox}
\textbf{NoSQL \& NewSQL Databases}

Designed for horizontal scalability, high throughput, loose schemas, and eventual consistency.

\smallskip
\begin{tabularx}{\linewidth}{l l l X}
  \toprule
  \textbf{Database} & \textbf{Model Type} & \textbf{Primary Lang.} & \textbf{Key Features \& Use Cases} \\
  \midrule
  Redis         & Key-Value         & C    & Disk-backed in-memory cache, fast Pub/Sub, key TTL expiration. \\
  Cassandra     & Column-Family     & Java & Massive dataset storage, CQL3 language. No JOINs, zero single point of failure. \\
  MongoDB       & Document          & C++  & JSON/BSON document store, built-in auto-sharding \& JS query language. \\
  ElasticSearch & Document / Search & Java & Distributed full-text search, JSON REST API, core component of ELK stack. \\
  \bottomrule
\end{tabularx}

\smallskip
\textbf{ELK Stack:} ElasticSearch (Search Engine) + Logstash (Log Processing) + Kibana (Visualization).
\end{conceptbox}

% =============================================================================
\section{Production Architecture Best Practices \& Reliability}

\begin{conceptbox}
\textbf{1.\ Decoupling}
\begin{itemize}
  \item Separate Web, Application, and Database services onto independent hardware nodes/containers.
  \item \textbf{Benefits:} Reduced blast radius (one failing service does not crash others), isolated resource allocations, easier troubleshooting, step towards microservices.
\end{itemize}
\end{conceptbox}

\begin{conceptbox}
\textbf{2.\ Scalability Principles}
\begin{itemize}
  \item \textbf{Statelessness:} The fundamental requirement for horizontal scaling. Application servers must NOT hold local session state.
  \item \textbf{Scale-Up (Vertical):} Upgrading to bigger/faster hardware (CPUs, RAM). Limited by physical limits \& high cost.
  \item \textbf{Scale-Out (Horizontal):} Adding more uniform nodes behind a load balancer. Highly scalable \& cost-effective.
  \item \textbf{Bottlenecks to monitor:} Compute (CPU/RAM) vs.\ I/O Performance (IOPS, Bandwidth, Latency).
\end{itemize}
\end{conceptbox}

\begin{conceptbox}
\textbf{3.\ Reliability Definitions \& Metrics}
\begin{itemize}
  \item \textbf{Durability:} Protection of data against loss/corruption (e.g., AWS S3 features 99.999999999\% / 11 9s durability).
  \item \textbf{Availability:} Percentage of total operational uptime (e.g., 99.99\% = 52.6 minutes downtime/year).
  \item \textbf{Reliability:} System capability to perform intended operations correctly while available.
  \item \textbf{Resiliency:} Ability of a system to self-heal and recover automatically from hardware/software failures.
  \item \textbf{Single Point of Failure (SPOF):} Any isolated component whose malfunction halts the entire system.
  \item \textbf{Chaos Engineering:} Testing system resiliency by deliberately injecting random failure (e.g., Netflix Chaos Monkey).
\end{itemize}
\end{conceptbox}

% =============================================================================
\section{Progressive Roadmap for Scaling Web Architecture}

Web architectures evolve incrementally from simple MVPs to multi-tiered resilient systems:

\begin{conceptbox}
\textbf{Stage 1: Single Server MVP}

\smallskip
Client Browser $\xrightarrow{\text{HTTP/HTTPS}}$ Single Server (Nginx + App + DB)

\smallskip
\textit{Evaluation:} Excellent for rapid prototype/demo; NOT production ready (SPOF, zero scaling).
\end{conceptbox}

\begin{conceptbox}
\textbf{Stage 2: Separate Database Server}

\smallskip
Client Browser $\to$ Web \& App Server $\xrightarrow{\text{Internal Network}}$ Dedicated DB Server

\smallskip
\textit{Evaluation:} Eliminates resource contention between CPU (App) and Disk I/O (DB). Removes DB from public internet (DMZ).
\end{conceptbox}

\begin{conceptbox}
\textbf{Stage 3: Load-Balanced Application Layer}

\smallskip
$\text{Client} \to \text{Load Balancer} \to
\begin{cases}
  \text{App Server 1} \\
  \text{App Server 2}
\end{cases}
\to \text{DB Server}$

\smallskip
\textit{Evaluation:} Introduces Horizontal Scale-Out. If App Server 1 crashes, Load Balancer routes traffic to App Server 2 without user interruption.
\end{conceptbox}

\begin{conceptbox}
\textbf{Stage 4: Adding HTTP Accelerators / Caching}

\smallskip
Client $\to$ HTTP Accelerator (Varnish/Nginx) $\to$ Load Balancer $\to$ App Servers $\to$ DB

\smallskip
\textit{Evaluation:} Caches static \& dynamic responses in memory. Drastically lowers response latency and reduces load on App Servers.
\end{conceptbox}

\begin{conceptbox}
\textbf{Stage 5: Master-Slave Database Replication}

\smallskip
$\text{App Servers} \xrightarrow{\text{Read/Write Split}}
\begin{cases}
  \text{Master DB (Writes)} \\
  \text{Slave DB 1 (Reads)} \\
  \text{Slave DB 2 (Reads)}
\end{cases}$

\smallskip
\textit{Evaluation:} Offloads heavy database read traffic onto read replicas (Slaves), leaving the Master node dedicated strictly to write transactions.
\end{conceptbox}

\begin{conceptbox}
\textbf{Stage 6: Fully Integrated High-Availability Enterprise Stack}

\smallskip
Client $\to$ Load-Balanced Caching Proxies $\to$ Load-Balanced App Servers $\to$ Master-Slave DB Cluster

\smallskip
\textit{Evaluation:} Maximizes availability, durability, throughput, and fault tolerance across all layers.
\end{conceptbox}

\vspace{14pt}
\textcolor{rulegray}{\rule{\linewidth}{1pt}}

\end{document}
